\documentclass[12pt]{article}
\usepackage[margin=1in]{geometry}
\usepackage{amsmath,amssymb,amsthm}

\newtheorem{theorem}{Theorem}
\newtheorem{lemma}[theorem]{Lemma}
\theoremstyle{definition}
\newtheorem{definition}[theorem]{Definition}

\newcommand{\R}{\mathbb{R}}
\newcommand{\eps}{\varepsilon}

\begin{document}

\section{Limits and continuity}

\begin{definition}
A function $f \colon \R \to \R$ is \emph{continuous at} $a$ if for every
$\eps > 0$ there exists $\delta > 0$ such that $|x - a| < \delta$ implies
$|f(x) - f(a)| < \eps$.
\end{definition}

\begin{lemma}\label{lem:sum}
If $f$ and $g$ are continuous at $a$, then so is $f + g$.
\end{lemma}

\begin{proof}
Let $\eps > 0$. Choose $\delta_1, \delta_2 > 0$ so that $|x - a| <
\delta_1$ implies $|f(x) - f(a)| < \eps/2$ and $|x - a| < \delta_2$ implies
$|g(x) - g(a)| < \eps/2$. Put $\delta = \min\{\delta_1, \delta_2\}$. Then
for $|x - a| < \delta$,
\begin{align*}
  \bigl|(f+g)(x) - (f+g)(a)\bigr|
    &\le |f(x) - f(a)| + |g(x) - g(a)| \\
    &< \frac{\eps}{2} + \frac{\eps}{2} = \eps. \qedhere
\end{align*}
\end{proof}

\begin{theorem}\label{thm:square}
The function $f(x) = x^2$ is continuous at every $a \in \R$.
\end{theorem}

\begin{proof}
Fix $a$ and $\eps > 0$, and let
\[
  \delta = \min\left\{ 1, \frac{\eps}{2|a| + 1} \right\}.
\]
If $|x - a| < \delta$ then $|x| < |a| + 1$, so $|x + a| < 2|a| + 1$ and
\[
  |x^2 - a^2| = |x - a|\,|x + a| < \delta\,(2|a| + 1) \le \eps. \qedhere
\]
\end{proof}

\begin{theorem}[Intermediate value theorem]
Let $f \colon [a,b] \to \R$ be continuous with $f(a) < 0 < f(b)$. Then
$f(c) = 0$ for some $c \in (a,b)$.
\end{theorem}

\begin{proof}
Let $S = \{ x \in [a,b] : f(x) < 0 \}$ and $c = \sup S$, which exists
because $S$ is nonempty and bounded above by $b$. We rule out the two
alternatives to $f(c) = 0$.

Suppose $f(c) < 0$. By continuity there is $\delta > 0$ with $f(x) < 0$ on
$(c - \delta, c + \delta) \cap [a,b]$, and $c < b$ since $f(b) > 0$; so some
point greater than $c$ lies in $S$, which is impossible.

Suppose $f(c) > 0$. Then $f > 0$ on some interval $(c - \delta, c]$, so
$c - \delta$ is an upper bound of $S$ smaller than $c$, which is again
impossible. Hence $f(c) = 0$, and $c \ne a, b$.
\end{proof}

\begin{theorem}
Let $f(x) = \begin{cases} x \sin(1/x) & x \neq 0, \\ 0 & x = 0. \end{cases}$
Then $f$ is continuous at $0$.
\end{theorem}

\begin{proof}
For $x \neq 0$ we have $|f(x) - f(0)| = |x|\,|\sin(1/x)| \le |x|$, so
$\delta = \eps$ works.
\end{proof}

\end{document}
